%==============================================================================
% MATH 5080
% Class:        14
% Date:         Oct 15th
%==============================================================================

\newpage

\section{Infinite Series}

\subsection{Basic properties of series}

\begin{definition}[Infinite Series]
    Given a sequence $\{ a_n \} \subset \mathbb{R}$.
    An infinite series is defined as
    \[
        \sum\limits_{k = 1}^{\infty}  a_k
    \]
\end{definition}

\begin{definition}[Partial sum]
    Given a sequence $\{ a_n \} \subset \mathbb{R}$.
    For all $n \geq 1$,
    the \textbf{$n$-th partial sum} is defined as 
    \[
        S_n \coloneq \sum\limits_{k=1}^{n} a_k
    \]
    The set $\{ S_n \}_{n \in \mathbb{N}} $ is called
    the \textbf{partial sums} of the series $\sum\limits_{n=1}^{\infty} a_n $.
\end{definition}


\begin{definition}[Convergence of a series]
    Given a sequence $\{ a_n \} \subset \mathbb{R}$,
    and partial sums $\{ S_n \}_n$,
    \begin{itemize}
        \item we say the series $\sum\limits_{k=1}^{\infty} a_k$ converges 
        to $S \in \mathbb{R}$
        if $ \{ S_n \} $ converges to $S \in \mathbb{R}$, this is denoted
        \[
            \sum\limits_{n = 1}^{\infty}  a_n = S
        \]
        \item we say the series diverges if $\{ S_n \} $ diverges.
    \end{itemize} 
\end{definition}

\begin{example}[Geometric Series]
    Fix $a \in \mathbb{R}$, its geometric series is 
    \[
    \sum\limits_{n = 1}^{\infty} a^{n-1} = 1 + a + a^{2} + \hdots
    \]
    This converges to $ \frac{1}{1-a}$ if $|a| < 1$, 
    and diverges otherwise.
    \medskip

    \textbf{Proof}: Recall
    \[
        \sum\limits_{k=1}^{n} a^{k-1} 
        = 1 + a + a^{2} + \hdots a^{n}
        = \frac{1-a^{n}}{1-a}
    \]
    Consider 3 cases
    \begin{itemize}
        \item $|a| < 1$: WTS $\lim_{n \to \infty}a^{n} = 0$
        Fix $\ep > 0$, take $N = \frac{\ln \ep}{\ln \left| a \right| }$.
        Then $\forall n > N$
        \begin{align*}
            n > \frac{\ln \ep}{\ln \left| a \right| } \\
            \implies &n \ln \left| a \right|  < \ln \ep
            \implies & \left| a \right|^{n} < \ep
            \implies & |a^{n} | < \ep
        \end{align*}
        and we are done.
        Hence 
        \[
            \sum\limits_{k = 1}^{\infty}  a^{k-1}
            = 
            \lim_{n \to \infty} \frac{1 - a^{n}}{1-a} 
            = \frac{1}{1-a}
        \]
        \item $|a| = 1$: $S_n = n$, diverges.
        \item $|a| > 1$: $\lim_{n \to \infty}a^{n}$ diverges.
    \end{itemize} 
\end{example}

\begin{theorem}[Series converges iff sequence of partial sums is Cauchy]
    $\sum\limits_{n = 1}^{\infty} a_n$ converges if and only if
    $\forall \ep >0, \exists N > 0$ s.t. 
    $m \geq n > N \implies \left| \sum\limits_{k = n}^{m} a_k \right| < \ep$.
\end{theorem}

\begin{proof}
    By definition and Cauchy Convergence Criterion,
    \begin{align*}
        & \sum\limits_{n = 1}^{\infty} a_n \text{ converges}  \\
        \iff &  \{ S_n \}  \text{ converges} \\
        \iff & \{ S_n \}  \text{ Cauchy} \\
        \iff & \forall \ep > 0, \exists N > 0 s.t. 
        \forall m > n > N, |S_n - S_m | = 
        \left| \sum\limits_{k = n + 1}^{m} a_k \right| 
        < \ep \\
        \iff & \forall \ep > 0, \exists N > 0 s.t. 
        \forall m \geq n > N_1, 
        \left| \sum\limits_{k = n}^{m} a_k \right| 
        < \ep \\
    \end{align*}
    The last equivalence (more specifically, its forward direction,
    requires a relabling, \remind{revisit this}
    ).
\end{proof}

\begin{corollary}[Convergence of series implies convergence of sequence to 0]
    If $\sum\limits_{n = 1}^{m} a_n$ converges
    then $\lim_{n \to \infty}a_n = 0$.
\end{corollary}

\begin{proof}
    By theorem, $\forall \ep > 0, \exists N > 0$ s.t. 
    $\forall m \geq n > N$
    \[
        \left| \sum\limits_{k=n}^{m} a_k  \right|  < \ep
    \]
    Take $m = n$, then $ \left| a_n  \right| < \ep$. 
    Hence $\lim_{n \to \infty}a_n = 0$.
\end{proof}

\begin{remark}
    \begin{itemize}
        \item The converse is not true. Convergence of a sequence
        to 0 does not imply convergence of series. For example, 
        harmonic series.
        \item The corollary is useful for showing a series diverges,
        by using its contrapositive.
    \end{itemize} 
\end{remark}

\begin{definition}[Absolute and conditional convergence]
    \hfill
    \begin{enumerate}
        \item We say $\sum\limits_{n = 1}^{\infty} a_n$ \textbf{converges
        absolutely} if $ \sum\limits_{n = 1}^{\infty} |a_n| $ converges.
        \item If $\sum\limits_{n=1}^{\infty} a_n$ converges 
        but $ \sum\limits_{n = 1}^{\infty} |a_n|$ diverges, then we say 
        $\sum\limits_{n = 1}^{\infty} a_n$ \textbf{converges conditionally}.
    \end{enumerate} 
\end{definition}

\begin{corollary}[Absolute convergence implies convergence]
    If $\sum\limits_{n = 1}^{\infty} a_n$ converges absolutely, 
    then $\sum\limits_{n = 1}^{\infty} a_n$ converges.
\end{corollary}

\begin{proof}
    Suffices to verify that the sequence of partial sums 
    is Cauchy. 
    \medskip

    Fix $\ep > 0$, since $\sum\limits_{}^{} |a_n|$ converges,
    $\exists N > 0$ s.t.
    \begin{align*}
        m \geq n > N 
        \implies &
        \left| \sum\limits_{k=n}^{m}|a_k|  \right|  < \ep \\
        \implies & \sum\limits_{k = n}^{m} |a_k| < \ep 
    \end{align*}
    By triangle inequality,
    \[
    \left| \sum\limits_{k = n}^{m} a_k \right|  
    \leq 
    \sum\limits_{k = n}^{m}  \left| a_k \right| 
    < \ep
    \]
\end{proof}

\begin{proposition}
    Let $\sum\limits_{n= 1}^{\infty} a_n$ and 
    $\sum\limits_{n=1}^{\infty} b_n$ converge,
    and $\alpha, \beta \in \mathbb{R}$,
    then 
    \[
        \sum\limits_{n= 1}^{\infty} \left( \alpha a_n + \beta b_n \right) 
        = \alpha \sum\limits_{n=1}^{\infty} a_n + \beta \sum\limits_{n =1}^{\infty} b_n
    \]
\end{proposition}

\begin{proof}
    Let $A_n = \sum\limits_{k= 1}^{n} a_k$,
    $B_n = \sum\limits_{k= 1}^{n} b_k$.
    By assumption
\end{proof}

\begin{remark}
    It is not generally true that 
    \[
        \sum\limits_{n = 1}^{\infty} (a_n b_n) 
        = 
        \left( \sum\limits_{n = 1}^{\infty} a_n \right) 
        \left( \sum\limits_{n = 1}^{\infty} b_n \right) 
    \]
\end{remark}

\subsection{Nonnegative series}

\begin{definition}[Nonnegative series]
    If $a_n \geq 0, \forall n \in \mathbb{N}$,
    then $\sum\limits_{n= 1}^{\infty} a_n$ is called a 
    \textbf{nonnegative series}.
\end{definition}

\begin{theorem}[Nonnegative series, theorem 1]
    \label{thm:nonnegative-series-1}
    Let $\sum\limits_{}^{} a_n$ be a nonnegative series, then 
    \begin{enumerate}
        \item $\sum\limits_{}^{} a_n $ converges $\iff$ 
        its partial sums $\{ S_n \} $ are bounded
        \item If $\sum\limits_{}^{} a_n$ diverges, then 
        \[
            \sum\limits_{n=1}^{\infty} a_n = \lim_{n \to \infty} S_n = \infty
        \]
    \end{enumerate} 
\end{theorem}

\begin{proof}
    $\{ S_n \} $ is a nonnegative increasing sequence. 
    \begin{enumerate}
        \item By MCT, since $ \{ S_n \} $ monotone,
        it converges if and only if it is bounded.
        \item $\sum\limits_{}^{} a_n$ diverges, then
        by (1), $\{ S_n \} $ unbounded from above.
        Hence $\forall M > 0, \exists N$ s.t. 
        $S_N > M$, 
        and $\forall n > N, S_n \geq S_N > M$.
        Hence 
        $\lim_{n \to \infty}S_n = \infty$.
    \end{enumerate} 
\end{proof}

\begin{theorem}[Nonnegative Series, Theorem 2]
    \label{thm:nonnegative-series-2}
    Given a series $\sum\limits_{}^{} a_n$ and let $\sum\limits_{}^{} b_k$
    be a series obtained from a regrouping of $\sum\limits_{}^{} a_n$ 
    (obtained by adding paranthesis), 
    then
    \begin{enumerate}
        \item $\sum\limits_{}^{} a_n$ converges $\implies$
        $\sum\limits_{}^{} b_k$ converges and 
        $\sum\limits_{}^{} a_n = \sum\limits_{}^{} b_k$
        \item If $\sum\limits_{}^{} a_n$ is nonnegative series, 
        then $\sum\limits_{}^{} a_n = A \iff \sum\limits_{}^{} b_k = A$.
    \end{enumerate} 
\end{theorem}

\begin{remark}
    When $a_n$ nonnegative, 
    convergence of series and convergence of 
    any regrouping is equivalent.
    \medskip

    For general $a_n$, we only have the forward direction.
    The converse (convergence of regrouping implying
    convergence of original series)
    is NOT true.
\end{remark}


\begin{proof}
    Let $S_n$ and $t_n$
    denote the $n$-th partial sum of $\sum\limits_{}^{} a_n$
    and $\sum\limits_{}^{} b_k$ respectively. Then
    $\{ t_n \} $ is a subsequence of $\{ S_n \} $.
    \medskip

    For the forward direction, which holds 
    for any convergent $\sum\limits_{}^{} a_n$
    \begin{align*}
        \sum\limits_{}^{} a_n \text{ converges} 
        &\implies \{ S_n \}  \text{ converges} \\
        & \implies \{ t_n \}  \text{ converges } \\
        & \implies \sum\limits_{}^{} b_k \text{ converges}
    \end{align*}
    
    For the reverse direction, $\sum\limits_{}^{} a_n$ nonnegative,
    $\sum\limits_{}^{} b_k$ converges.
    Since $S_n$ is increasing and has a convergent subsequence
    $t_n$, therefore $S_n$ must converge to the same limit.
    \begin{itemize}
        \item $t_n$ converges implies $t_n$ bounded.
        \item $S_n$ monotone, and has a bounded subsequence $t_n$,
        hence $S_n$ is bounded.
        \item By MCT, $S_n$ converges.
    \end{itemize} 
\end{proof}

\begin{corollary}
    Let $\{ S_n \} $ be the sequence of 
    partial sums of nonnegative series $\sum\limits_{}^{} a_n $.
    If $\{ S_n \} $ has a bounded subsequence (or equivalently convergent)
    then $\sum\limits_{}^{} a_n $ converges.
\end{corollary}

\subsection{Convergence Rules for Nonnegative series}

\begin{theorem}[Comparison test]
    \label{thm:nonnegative-series-3}
    Let $\sum\limits_{}^{} a_n, \sum\limits_{}^{} b_n$ be 
    two nonnegative series. If $\exists N > 0$ s.t. 
    $\forall n > N, a_n \leq b_n$, then
    \begin{enumerate}
        \item $\sum\limits_{}^{} b_n$ converges $\implies$ $\sum\limits_{}^{} a_n$ converges.
        \item $\sum\limits_{}^{} a_n$ diverges $\implies$ $\sum\limits_{}^{} b_n$ diverges.
    \end{enumerate} 
\end{theorem}

\begin{proof}
    WLOG, assume $a_n \leq b_n, \forall n \in \mathbb{N}$. \remind{
        Check why we can do this.
    }
    
    \textbf{Proof of (1)}:
    Let $S_n^{a}, S_n^{b}$ denote the $n$-th partial sums of 
    $\sum\limits_{}^{} a_n, \sum\limits_{}^{} b_n$.
    \begin{align*}
        & \sum\limits_{}^{} b_n \text{ converges} \\
        \implies & S_n^{b} \text{ bounded} \\
        \implies & \exists M > 0 \text{ s.t. } |S_n^{b} | \leq M, \forall n \in \mathbb{N} \\
        \implies & |S_n^{a} |\leq \left| S_n^{b} \right|  \leq M, \forall n \in \mathbb{N} \\
        \implies & S_n^{a} \text{ bounded} \\
        \implies & \sum\limits_{}^{} a_n \text{ converges}
    \end{align*}
    \textbf{Proof of (2)}:
    \begin{align*}
        & \sum\limits_{}^{} a_n \text{ diverges} \\
        \implies & S_n^{a} \text{ unbounded} \\
        \implies & S_n^{b} \text{ unbounded} \\
        \implies & \sum\limits_{}^{} b_n \text{ diverges}
    \end{align*}
\end{proof}

\begin{example}[$p$-series]
    $\sum\limits_{n=1}^{\infty}  \frac{1}{n^{p}}$
    converges if $p > 1$, diverges if $ p \leq 1$. 
    \medskip 

    \textbf{Proof}: 
    \medskip

    \textbf{Step 1}: $p > 1$,
    suffices to show a regrouping converges.
    Consider 
    \[
        S_{2^{k} - 1} = \frac{1}{1} + \frac{1}{2^{p}}
        + \frac{1}{3^{p}} 
        + \frac{1}{(2^{k} - 1)p}
    \]
    Let $b_k = S_{2^{k} - 1} - S_{2^{k-1}-1}, \forall k \geq 2$, and $b_1 = a_1$.
    i.e. 
    \begin{align*}
        b_1 &= a_1 \\
        b_2 &= S_3 - S_1 = a_2 + a_3 \\
        b_3 &= S_7 - S_3 = a_4 + a_5 + a_6 + a_7
    \end{align*}
    WTS $\sum\limits_{}^{} b_k $ converges.
    We use comparison test with geometric series.
    \begin{align*}
        b_k &= S_{2^k- 1} - S_{2^{k-1} - 1} \\
        &= \sum\limits_{n = 2^{k-1}}^{2^k}  \frac{1}{n^{p}} \\
        & \leq \frac{1}{(2^{k-1})^p} (2^{k} - 2^{k-1} ) \\
        & =\frac{2^{k-1}}{(2^{k-1})^p} \\
        &= (2^{k-1})^{1-p}  \\
        &= \left( \frac{1}{2^{p-1}} \right)^{k-1}
    \end{align*}

    From geometric series, we know that 
    \[
        \sum\limits_{k=1}^{\infty}  \left( 
            \frac{1}{2^{P - 1}}
        \right)^{k-1}
    \]
    converges, hence $\sum\limits_{}^{} b_k$ converges.
    \medskip

    \textbf{Step 2}: $p \leq 1$: since $ \frac{1}{n^{p}} \geq \frac{1}{n}$,
    by comparison test, $ \sum \frac{1}{n^{p}}$ diverges.
\end{example}

\begin{theorem}[Ratio test]
    \label{thm:nonnegative-series-4}
    Let $\sum\limits_{}^{} a_n$ be a nonnegative series, $a_n \neq 0$. Then 
    \begin{enumerate}
        \item If $\limsup_{n \to \infty} \frac{ a_{n + 1}}{a_n} < 1$, then 
        $\sum\limits_{}^{} a_n$ converges.
        \item If $\liminf_{n \to \infty} \frac{a_{n + 1}}{a_n} > 1$, then 
        $\sum\limits_{}^{} a_n$ diverges.
    \end{enumerate} 
\end{theorem}

\begin{remark}
    If the limsup or liminf is 1, then it is inconclusive.
\end{remark}

\begin{example}
    $\sum\limits_{n=1}^{\infty} \frac{n}{2^{n}}$ converges,
    by ratio test.
    \begin{align*}
        \frac{a_{n + 1}}{a_n} = \frac{n+1}{2^{n +1}} \frac{2^{n}}{n}
        = \frac{1}{2} \frac{n+ 1}{n} 
        \to \frac{1}{2}
    \end{align*}
\end{example}

\begin{proof}
    \begin{enumerate}
        \item Let 
        $r \in (\limsup_{n \to \infty} \frac{ a_{n + 1}}{a_n}, 1)$,
        then $\exists N > 0$ s.t. $\forall n > N$, 
        $ \frac{a_{n + 1}}{a_n} < r$.
        Hence 
        \begin{align*}
            a_{n + 1}
            &= 
            \frac{a_{n + 1}}{a_n}
            \frac{a_{n}}{a_{n-1}}
            \frac{a_{n -1 }}{a_{n-2}} 
            \hdots
            \frac{a_{N +2 }}{a_{N+1}} 
            a_{N + 1} \\
            & \leq r^{n - N} a_{N + 1} \\
            & = r^n \frac{a_{N + 1}}{r^N} \\
            &= cr^n \text{ for some constant } c
        \end{align*}
        Since $r < 1 \implies \sum\limits_{}^{} cr^n$
        converges, by comparison test, 
        $\sum\limits_{}^{} a_n$ converges
        \remind{revisit this}
        \item \remind{fill in mirrored proof}
    \end{enumerate} 
\end{proof}


\begin{theorem}[Root test]
    \label{thm:nonnegative-series-5}
    Let $\sum\limits_{}^{} a_n$ be a nonnegative series, then 
    \begin{enumerate}
        \item If $\limsup_{n \to \infty} (a_n)^{ \frac{1}{n}} < 1$ then
        $\sum\limits_{}^{} a_n$ converges.
        \item If $\limsup_{n \to \infty} (a_n)^{ \frac{1}{n}} > 1$ then
        $\sum\limits_{}^{} a_n$ diverges.
    \end{enumerate} 
\end{theorem}

\begin{proof}
    \begin{enumerate}
        \item Let $r \in \left( \limsup_{n \to \infty} a_n^{ \frac{1}{n}}, 1\right) $
        then $\exists N$ s.t $\forall n > N$,
        $a_n^{1/n} < r \implies a_n < r^{n}$. 
        Since $r < 1$, $\sum\limits_{}^{} r^n$ converges.
        Hence by comparison test, $\sum\limits_{}^{} a_n$ converges
        \remind{revisit this}
        \item \remind{fill in mirrored proof}
    \end{enumerate} 
\end{proof}

\begin{example}[Example where ratio test fails but root test works]
    Consider the series
    \[
        \frac{1}{2} + \frac{1}{3} + \frac{1}{2^{2}} + \frac{1}{3^{2}}
        + \frac{1}{2^{3}} + \frac{1}{3^{3}} + \hdots
    \]
    that is, $a_{2k-1} = \frac{1}{2^{k}}$ and $a_{2k} = \frac{1}{3^{k}}$
    for $k \geq 1$.
    \medskip

    \textbf{Ratio test}: the ratios alternate between two behaviours.
    \begin{align*}
        \frac{a_{2k}}{a_{2k-1}} &= \frac{1/3^{k}}{1/2^{k}}
        = \left( \frac{2}{3} \right)^{k} \to 0 \\
        \frac{a_{2k+1}}{a_{2k}} &= \frac{1/2^{k+1}}{1/3^{k}}
        = \frac{1}{2} \left( \frac{3}{2} \right)^{k} \to \infty
    \end{align*}
    Hence
    $\liminf_{n \to \infty} \frac{a_{n+1}}{a_n} = 0$ and
    $\limsup_{n \to \infty} \frac{a_{n+1}}{a_n} = \infty$.
    Neither $\limsup < 1$ nor $\liminf > 1$ holds,
    so the ratio test does not apply.
    \medskip

    \textbf{Root test}: $\limsup$ is the largest subsequential limit, 
    so we compute the limit of $\sqrt[n]{a_n}$ along the odd and the even
    terms separately.
    \begin{itemize}
        \item $n = 2k-1$ (odd), $a_n = \frac{1}{2^{k}}$:
        \[
            \sqrt[n]{a_n} = \left( \frac{1}{2^{k}} \right)^{\frac{1}{2k-1}}
            = 2^{-\frac{k}{2k-1}}
            = 2^{-\frac{1}{2 - 1/k}}
            \xrightarrow{k \to \infty} 2^{-1/2} = \frac{1}{\sqrt{2}}
        \]
        \item $n = 2k$ (even), $a_n = \frac{1}{3^{k}}$:
        \[
            \sqrt[n]{a_n} = \left( \frac{1}{3^{k}} \right)^{\frac{1}{2k}}
            = 3^{-1/2} = \frac{1}{\sqrt{3}}
        \]
    \end{itemize}
    These are the only subsequential limits, so
    \[
        \limsup_{n \to \infty} \sqrt[n]{a_n}
        = \max \left\{ \frac{1}{\sqrt{2}}, \frac{1}{\sqrt{3}} \right\}
        = \frac{1}{\sqrt{2}} < 1
    \]
    By the root test, the series converges.
\end{example}


\subsection{Discussion of $e$}

\begin{definition}[$e$]
    $e$ is defined as the sum of the infinite series
    \[
        e \coloneq \sum\limits_{n = 0}^{\infty} \frac{1}{n!}
        = 1 + \frac{1}{1!} + \frac{1}{2!} + \hdots
    \]
\end{definition}

\remind{Add in discussion of $e$}.
